\documentclass{package/notes}

%%%%%%%%%% Default Package %%%%%%%%%%%%%
\usepackage{package/color-env}
%%%%%%%%%% %%%%%%%%%%%%%%% %%%%%%%%%%%%%

%%%%%%%%%% Required Packages %%%%%%%%%%%%%
\usepackage{background}
\usepackage[object=vectorian]{pgfornament} %% used in title.tex
\usepackage{calligra} %%% (optional) to make the Title text beautiful 

%%%%%% Optional Packages %%%%%%%
\usepackage{lettrine} %% for nice looking 
\usepackage{GoudyIn} %% first Letter of the paragraph
%%%%%%%%%%%%%%%%%%%%%%%%

\usepackage{lipsum}  %% for dummy text 
\usepackage{amssymb,amsmath,amsfonts}  %%% for maths
%%%%%%%%%%%%%%%%%%%%%%%%%%%%%%%%%%%%%

%%%%%%%%%%%%%%%%%%%%%%%%%%%%%%%%%%%%%
\renewcommand{\LettrineFontHook}{\color{black}\GoudyInfamily{}}
\LettrineTextFont{\itshape}
\setcounter{DefaultLines}{3}%
%%%%%%%%%%%%%%%%%%%%%%%%%%%%%%%%%%%%%

%%%% Bibliography %%%%%%%%%
% Required packages are included in notes class
% Can be tweaked in the notes.cls file itself

\addbibresource{resource/references.bib}

\begin{document}

\begin{titlepage}

%%%%%%%%%%%%%%%%%%%%%%%%%%%%%%%%%%%% Inspired From  %%%%%%%%%%%%%%%%%%%%%%%%%%%%%%%%%%%%%%%%%%%%%%%%% 
%%%%  https://www.reddit.com/r/LaTeX/comments/j9d739/hello_world_in_latex_is_a_lot_cooler/   %%%%%
%%%%%%%%%%%%%%%% If this doesn't look nice then you may remove it %%%%%%%%%%%%%%%%%%%%%%%%%%

\backgroundsetup{
scale=1,
opacity=1,
angle=0,
color=black,
contents={
\begin{tikzpicture}[color=black, every node/.style={inner sep= 15pt}]
\node (NW) [anchor=north west] at (current page.north west){\pgfornament[width=2.5cm] {131}};
\node (NE) [anchor=north east] at (current page.north east){\pgfornament[width=2.5cm, symmetry=v]{131}};
\node (SW) [anchor=south west] at (current page.south west){\pgfornament[width=2.5cm, symmetry=h]{131}};
\node (SE) [anchor=south east] at (current page.south east){\pgfornament[width=2.5cm, symmetry=c]{131}};
\foreach \i in {-4,0,4}
\node[anchor=north,xshift=\i cm] at (current page.north){\pgfornament[scale=0.25,symmetry=v]{71}};
\foreach \i in {-4,0,4}
\node[xshift=\i cm, yshift=32.25 pt] at (current page.south){\pgfornament[scale=0.25,symmetry=v]{71}};
\foreach \i in {-8,-4,0,4,8}
\node[yshift=\i cm, xshift=32.25pt, rotate=90] at (current page.west){\pgfornament[scale=0.25,symmetry=v]{71}};
\foreach \i in {-8,-4,0,4,8}
\node[yshift=\i cm, xshift=-32.25pt, rotate=90] at (current page.east){\pgfornament[scale=0.25,symmetry=v]{71}};
\foreach \i in {-11,-9,...,7,9}
\node[anchor=west, yshift=\i cm, xshift=52.25pt, rotate=90] at (current page.west){\pgfornament[scale=0.1]{80}};
\foreach \i in {-11,-9,...,7,9}
\node[anchor=east, yshift=\i cm, xshift=-52.25pt, rotate=-90] at (current page.east){\pgfornament[scale=0.1]{80}};
\end{tikzpicture}
}}

%%%%%%%%%%%%%%%%%%%%%%%%%%%%%%%%%%%%%%%%%%%%%%%%%%%%%%%%%%%%%%%%%%%%%%%%%%%%%%%%%%%%%%%%%%%%%%%%%%%%%%%%%%%%%%%%%%

\centering % Centre everything on the title page
		
\scshape % Use small caps for all text on the title page

\vspace*{\baselineskip} % White space at the top of the page

%------------------------------------------------
%	Title
%------------------------------------------------

\rule{\textwidth}{1.6pt}\vspace*{-\baselineskip}\vspace*{2pt} % Thick horizontal rule
\rule{\textwidth}{0.4pt} % Thin horizontal rule

\vspace{0.75\baselineskip} % Whitespace above the title

{\huge \calligra{ The Title }\\} % Title

\vspace{0.75\baselineskip} % Whitespace below the title

\rule{\textwidth}{0.4pt}\vspace*{-\baselineskip}\vspace{3.2pt} % Thin horizontal rule
\rule{\textwidth}{1.6pt} % Thick horizontal rule

\vspace{2\baselineskip} % Whitespace after the title block

%------------------------------------------------
%	Subtitle
%------------------------------------------------

\LARGE{Course Name} 

\vspace*{3\baselineskip} % Whitespace under the subtitle



\vspace{0.5\baselineskip} 

{\scshape   \LARGE Professor\\ } % Editor list

\vspace{0.2\baselineskip} 

\textit{\Large University} 

\vfill 

%------------------------------------------------
% Author
%------------------------------------------------

\begin{figure}[!h]
    \centering
    \includegraphics[width = 3cm, height= 3cm]{resource/icon.jpg}%% include the university icon here
\end{figure}
\vspace{0.3\baselineskip} 


{\large Edited by\\  Your name}
\end{titlepage}

\newpage

\backgroundsetup{contents={}} %% to remove background and watermark from other pages
\tableofcontents

%\newpage
\chapter{A Nice Title}

\lettrine{L}{o} %% for the calligraphic First Letter
rem ipsum dolor sit amit consectetuer adipiscingelit \footnote[1]{Here is an example footnote. \LaTeX{} can be customized in many different ways}.  
\lipsum[1][2-100] \cite{wiki:latex}



\section{Here goes the section!}
\subsection{And the subsection}

\begin{definition}[DEF NAME]{def:label} %%  [can be kept empty]
	\lipsum[1][1-3] %% for dummy text
	Here is the nice looking Schrodinger's Equation \cite{dirac}
	\[i\hbar\frac{\partial}{\partial t} \varPsi(\mathbf{r},t) = \left [ \frac{-\hbar^2}{2m}\nabla^2 + V(\mathbf{r},t)\right ] \varPsi(\mathbf{r},t)\]
\end{definition}

\begin{theorem}[THM NAME]{thm:label}%%  [can be kept empty]
	\lipsum[2][1-5] %% for dummy text
\end{theorem}

\begin{proposition}[PROP NAME]{prop:label}%%  [can be kept empty]
	\lipsum[4][1-4] %% for dummy text
\end{proposition}

\begin{lemma}[LEM NAME]{lem:label}%%  [can be kept empty]
	\lipsum[3][1-3] %% for dummy text
\end{lemma}

\begin{corollary}[COR NAME]{cor:label}%%  [can be kept empty]
	\lipsum[5][1-3] %% for dummy text
\end{corollary}

\begin{problem}%%  [can be kept empty]
	\lipsum[1][1-3] %% for dummy text
\end{problem}
\begin{proof}[Proof Head]
	\lipsum[1][1-10] %% for dummy text
\end{proof}
\nocite{knuthwebsite}

\section{Mathematics and Images}
\begin{figure}[!h]
	\centering
	\includegraphics[height = 0.5\textwidth]{resource/ctan_lion.png}
	\caption{We can add images too!! Isn't that good.}
	\label{img:anything}
\end{figure}

\lettrine{M}{a} xwell's equations are a set of coupled partial differential equations that, together with the Lorentz force law, form the foundation of classical electromagnetism, 
classical optics, and electric circuits. The equations provide a mathematical model for electric, optical, and radio technologies, such as power generation, electric
 motors, wireless communication, lenses, radar etc. They describe how electric and magnetic fields are generated by charges, currents, and changes of the fields.
  The equations are named after the physicist and mathematician James Clerk Maxwell, who, in 1861 and 1862, published an early form of the equations that included 
  the Lorentz force law. Maxwell first used the equations to propose that light is an electromagnetic phenomenon. 

  Here are the Maxwell's equations in the differential form

\[
\begin{gathered}
	\nabla \cdot \mathbf {E} ={\frac {\rho }{\varepsilon _{0}}} \\
	 \nabla \cdot \mathbf {B} ={0} \\ \tag*{(1)}
	\nabla \times \mathbf {E} =-{\frac {\partial \mathbf {B} }{\partial t}}\\
	 \nabla \times \mathbf {B} =\mu _{0}\left(\mathbf {J} +\varepsilon _{0}{\frac {\partial \mathbf {E} }{\partial t}}\right)
\end{gathered}
\]





\pagebreak

\medskip

\printbibliography[heading=bibintoc,title={\centering Bibliography}]

\end{document}
